\section{Consequences and comparison}

Theorems~\ref{thm:multiset} and \ref{thm:distinct} imply, for example,
\[
\limsup_{T\to\infty}
\frac{N(T,2T)-M_{\leq3}(T,2T)}{N(T,2T)}
\leq0.112379991826645\ldots,
\]
so at most $11.2380\%$ of the zero multiset may be supported on locations of multiplicity at least four.  Among distinct locations, at most $3.0681\%$ may have multiplicity at least four.

For $j\geq2$, define the exact-multiplicity multiset and distinct counts
\[
E_j(T_1,T_2):=\sum_{\substack{T_1<\gamma\leq T_2\\m_\rho=j}}m_\rho,
\qquad
E_{d,j}(T_1,T_2):=\#\{\rho:T_1<\gamma\leq T_2,\ m_\rho=j\}.
\]
Taking $r=j-1$ in the two profile theorems gives the following exact-multiplicity upper bounds.

\begin{corollary}[Exact multiplicity profile]\label{cor:exact}
One has
\begin{align*}
\limsup_{T\to\infty}\frac{E_2(T,2T)}{N(T,2T)}
&\leq \kMT-1,\\
\limsup_{T\to\infty}\frac{E_3(T,2T)}{N(T,2T)}
&\leq \frac{\kMT-1}{2},\\
\limsup_{T\to\infty}\frac{E_j(T,2T)}{N(T,2T)}
&\leq B\frac{j}{j-2}\qquad(j\geq4),
\end{align*}
and
\begin{align*}
\limsup_{T\to\infty}\frac{E_{d,2}(T,2T)}{N_d(T,2T)}
&\leq \frac{\kMT-1}{3-\kMT},\\
\limsup_{T\to\infty}\frac{E_{d,3}(T,2T)}{N_d(T,2T)}
&\leq \frac{\kMT-1}{2(4-\kMT)},\\
\limsup_{T\to\infty}\frac{E_{d,j}(T,2T)}{N_d(T,2T)}
&\leq \frac{B}{j-2-(j-1)B}\qquad(j\geq4).
\end{align*}
Every bound is sharp within the scalar information model of Theorem~\ref{thm:sharp}.
\end{corollary}

\begin{proof}
The exact-multiplicity-$j$ set is contained in the set of locations of multiplicity at least $j$, so the inequalities follow from Theorems~\ref{thm:multiset} and \ref{thm:distinct} with $r=j-1$.  The sharpness distributions for $r=j-1$ place all of their high-multiplicity mass at multiplicity exactly $j$.
\end{proof}

For fixed $r$, these are unconditional two-trace bounds.  They are weaker numerically than the GRH-dependent higher-correlation results of \cite{GDL2025}; for instance, that work obtains $0.9614$ at multiplicity at most two and $0.9787$ at multiplicity at most three.  The distinction is that no RH or higher-level correlation input is used here.

Corollary~\ref{cor:exact} also complements the unconditional exact-multiplicity theorem of \cite{STTB2022}.  Their estimate is uniform in $j$ and decays exponentially, becoming effective relative to elementary simple-zero bounds for extremely large multiplicity.  The present estimates are fixed-$j$ consequences of two traces and are aimed at the low-multiplicity regime; in particular, the multiset upper bound $Bj/(j-2)$ is not asserted uniformly when $j$ grows with $T$.

As $r\to\infty$, the multiset lower bound tends to
\[
1-B=0.943810004086677\ldots,
\]
whereas the distinct-location lower bound tends to one with tail
\[
1-\frac{D_{\leq r}}{N_d}
\leq \frac{B}{r}+O\!\left(\frac1{r^2}\right).
\]
Thus the scalar two-trace information forces high multiplicities to occupy a vanishing fraction of distinct zero locations, although it still permits a fixed fraction of the zero multiset to escape to multiplicities growing with height.

\paragraph{Separate derivative problem.}
The analogous optimization for zeros of $\xi'$ is not used in this paper.  It requires profile-specific analytic integration beyond the scalar bookkeeping above and is therefore kept in a separate research package rather than folded into the present theorem claim.

\section{Exact verification, source audit, and formalization path}\label{sec:verification}

The accompanying standard-library Python certificate uses exact rational arithmetic.  The identity
\[
\kMT=\frac12+\frac{\cos(1/\sqrt2)}{\sin(1/\sqrt2)/(1/\sqrt2)}
\]
allows both numerator and denominator to be enclosed by alternating Taylor series with rational terms:
\[
\cos(1/\sqrt2)=\sum_{n\geq0}\frac{(-1)^n}{2^n(2n)!},
\qquad
\frac{\sin(1/\sqrt2)}{1/\sqrt2}
=\sum_{n\geq0}\frac{(-1)^n}{2^n(2n+1)!}.
\]
The certificate encloses $\kMT$, $B$, every table entry, and every exact-multiplicity upper bound by rational intervals.  All printed lower bounds are rounded downward; all printed upper bounds are rounded upward.  A second script uses exact symbolic algebra to verify the parameter simplifications, the square identities in the sharpness proof, and the pointwise majorants over their complete integer ranges.

The source audit at commit \path{3635e74826a4c1fcece7d1cd2b6fa75e43a00510} records the following dependency chain:
\Needspace{12\baselineskip}
\begin{description}[leftmargin=1.3em,style=nextline,font=\normalfont\ttfamily]
\item[ZeroSide.Mult block normalization and Poisson lemmas]
identify the normalized zero compression and verify the hypotheses of the generic matrix theorem;
\item[ZeroSide.RankTraceMult.rank\_trace\_mult\_k\_le]
arbitrary-$c$ inequality \eqref{eq:ranktracegeneric};
\item[Assembly.ctr\_sub\_frobSq\_perturb]
arbitrary trace coefficient in \eqref{eq:perturb};
\item[ThmD.tracesBoundsD\_concrete and tendsto\_cRatio\_concrete]
fixed-$\lambda$ Montgomery--Taylor trace asymptotics;
\item[eventually\_tailPackageD and local-count inputs]
$B_T\to0$ and $O(T^{1/2}\log T)$ boundary multiplicity.
\end{description}

A complete Lean extension now has a finite scope:
\begin{enumerate}[leftmargin=2em]
\item define $M_{\leq r}$, $D_{\leq r}$, and the exact-multiplicity counts;
\item package Proposition~\ref{prop:master} by combining the existing generic rank--trace and perturbation declarations with the existing Theorem-D trace/tail endgame;
\item formalize the elementary scalar majorants and sharpness identities;
\item instantiate at $c=2$, $c=3$, and $c=2+\sqrt2$, then reuse the existing dyadic and fixed-Dirichlet-$L$ summation layers;
\item run the pinned complete build and the repository's \texttt{\#print axioms} audit.
\end{enumerate}
No new explicit-formula, prime-side, or zero-localization theorem is needed.  Lean and Lake are not installed in the present runtime, so the extension supplied with this draft is a source-level blueprint rather than a compiled formal proof.
